\begin{lemma}[Claim T$'$: the supersolution bound along a stopped lattice portion]
Let $(\Omega,\mathcal M,\mu)$ be a probability space, let
$(\mathcal G_n)_{n\in\N}$ be a filtration of $\mathcal M$, let
$\xi:\N\times\Omega\to\R$, let $\Xi:\N\times\Omega\to\Z$, let
$S:\N\times\Omega\to\Z$, let $\tau,\sigma,\varsigma:\Omega\to\N$ and let
$a,b,q\in\N$. Assume:
\begin{enumerate}
\item\label{it:slsb-pm} $\xi_n(\omega)=1$ or $\xi_n(\omega)=-1$ for all
$n,\omega$;
\item\label{it:slsb-ximeas} $\xi_n$ is $\mathcal G_{n+1}$-measurable for every
$n$;
\item\label{it:slsb-xicond} $\mathbb E[\xi_n\mid\mathcal G_n]=0$ $\mu$-almost
everywhere, for every $n$;
\item\label{it:slsb-Xi} $\Xi_n(\omega)=1$ if $\xi_n(\omega)=1$ and
$\Xi_n(\omega)=-1$ otherwise;
\item\label{it:slsb-step} $S_{n+1}=S_n+\Xi_n$ on $\{n<\tau\}$ and
$S_{n+1}=S_n$ on $\{n\ge\tau\}$;
\item\label{it:slsb-range} $a\le S_n(\omega)\le q$ for all $n,\omega$;
\item\label{it:slsb-st} $\sigma$ and $\varsigma$ are $(\mathcal G_n)$-stopping
times and both are $\mathcal M$-measurable as maps $\Omega\to\N$;
\item\label{it:slsb-order} $\sigma\le\varsigma$ and $\varsigma\le\tau$
pointwise;
\item\label{it:slsb-stG} for every $(\mathcal G_n)$-stopping time
$\zeta:\Omega\to\N$ and every $n$, $\{\zeta\le n\}\in\mathcal G_n$;
\item\label{it:slsb-idx} if $\zeta,\eta:\Omega\to\N$ are measurable and
$G_{n,m}:\Omega\to\R$ is measurable for all $n,m$, then
$\omega\mapsto G_{\zeta(\omega),\eta(\omega)}(\omega)$ is measurable;
\item\label{it:slsb-bdd} every $g:\Z\to\R$ is bounded on
$\{i\in\Z:a\le i\le q\}$;
\item\label{it:slsb-meas} for every $g:\Z\to\R$ and every $k\le n$, the
function $\omega\mapsto g(S_k(\omega))$ is $\mathcal G_n$-measurable, and for
every $g$ and $k$ it is $\mathcal M$-measurable;
\item\label{it:slsb-win} $a<S_n(\omega)<b$ whenever
$\sigma(\omega)\le n<\varsigma(\omega)$.
\end{enumerate}
Let $h,u:\Z\to\R$ and $A\subseteq\Omega$ satisfy $u(i)\ge0$ for every
$i\in\Z$, the discrete supersolution inequality
\begin{equation}\label{eq:slsb-super}
  \frac{h(i+1)+h(i-1)}{2}-h(i)\le-u(i)\qquad(a<i<b),
\end{equation}
$A\in\mathcal M$, and $A\cap\{\sigma\le n\}\in\mathcal G_n$ for every $n$. Then
for every $n\in\N$,
\begin{equation}\label{eq:slsb-bound}
  \int\mathbf1_A(\omega)
     \Bigl(h\bigl(S_{(\sigma(\omega)\vee n)\wedge\varsigma(\omega)}(\omega)\bigr)
       +\sum_{\sigma(\omega)\le m<n\wedge\varsigma(\omega)}
          u\bigl(S_m(\omega)\bigr)\Bigr)d\mu(\omega)
  \le\int\mathbf1_A(\omega)\,h\bigl(S_{\sigma(\omega)}(\omega)\bigr)d\mu(\omega).
\end{equation}
\end{lemma}
\begin{proof}
By \ref{it:slsb-bdd} choose $C_h$, $C_h^+$, $C_h^-$ and $C_u$ bounding
$|h(i)|$, $|h(i+1)|$, $|h(i-1)|$ and $|u(i)|$ respectively for $a\le i\le q$;
by \ref{it:slsb-range} every value $S_m(\omega)$ lies in that range, so all four
bounds apply at every $S_m(\omega)$. Write $\iota=\mathbf1_A$, which is
measurable because $A\in\mathcal M$, satisfies $0\le\iota\le1$ and in particular
$|\iota|\le1$. For $n\in\N$ set
\[
  Y_n(\omega)=\iota(\omega)\Bigl(
    h\bigl(S_{(\sigma(\omega)\vee n)\wedge\varsigma(\omega)}(\omega)\bigr)
    +\sum_{\sigma(\omega)\le m<n\wedge\varsigma(\omega)}u\bigl(S_m(\omega)\bigr)
  \Bigr),
\]
so that \eqref{eq:slsb-bound} asserts $\int Y_n\,d\mu\le\int\iota\,h(S_\sigma)\,d\mu$.

\emph{Each $Y_n$ is integrable.} The clipped index
$(\sigma\vee n)\wedge\varsigma$ is a measurable $\N$-valued function of $\omega$
built from the measurable $\sigma,\varsigma$ of \ref{it:slsb-st}, so
$\omega\mapsto h(S_{(\sigma\vee n)\wedge\varsigma(\omega)}(\omega))$ is
measurable by \ref{it:slsb-idx} applied to the doubly-indexed measurable family
$(i,m,\omega)\mapsto h(S_{(i\vee n)\wedge m}(\omega))$, each member of which is
measurable by \ref{it:slsb-meas}. The sum is a finite sum, over the
$\sigma,\varsigma$-determined range, of the measurable functions
$\omega\mapsto u(S_m(\omega))$, hence measurable by the same two hypotheses.
So $Y_n$ is measurable; it is bounded by $C_h+n\,C_u$, since the sum has at
most $n$ terms and each is at most $C_u$ in absolute value; and a bounded
measurable function is integrable because $\mu$ is a probability measure.

\emph{The value at $n=0$.} Here $(\sigma\vee0)\wedge\varsigma
=\sigma\wedge\varsigma=\sigma$ by \ref{it:slsb-order}, and the summation range
$[\sigma,0\wedge\varsigma)=[\sigma,0)$ is empty. Hence
\begin{equation}\label{eq:slsb-Y0}
  \int Y_0\,d\mu=\int\iota\,h(S_\sigma)\,d\mu .
\end{equation}

\emph{The one-step identity.} Fix $n$. We claim that pointwise
\begin{equation}\label{eq:slsb-split}
  Y_{n+1}-Y_n
   =\iota\cdot D_n+\iota\cdot c_n\,\xi_n ,
\end{equation}
where, writing $P_n=\{\sigma\le n<\varsigma\}$,
\[
  D_n=\mathbf1_{P_n}\Bigl(\frac{h(S_n+1)+h(S_n-1)}{2}-h(S_n)+u(S_n)\Bigr),
  \qquad
  c_n=\mathbf1_{P_n}\,\frac{h(S_n+1)-h(S_n-1)}{2}.
\]
We check this in three cases, which exhaust $\Omega$.

On $\{\sigma>n\}$ we have $\sigma\ge n+1$, so
$(\sigma\vee n)\wedge\varsigma=\sigma\wedge\varsigma=\sigma$ and
$(\sigma\vee(n+1))\wedge\varsigma=\sigma\wedge\varsigma=\sigma$, using
$\sigma\le\varsigma$; and both summation ranges are empty because
$n\wedge\varsigma\le n<\sigma$ and $(n+1)\wedge\varsigma\le n+1\le\sigma$.
Hence $Y_{n+1}=Y_n$, and $P_n$ does not occur, so both terms on the right of
\eqref{eq:slsb-split} vanish.

On $\{\sigma\le n\}\cap\{\varsigma\le n\}$ we have
$(\sigma\vee n)\wedge\varsigma=n\wedge\varsigma=\varsigma$ and likewise
$(\sigma\vee(n+1))\wedge\varsigma=\varsigma$, and both summation ranges equal
$[\sigma,\varsigma)$. Hence $Y_{n+1}=Y_n$, and again $P_n$ does not occur.

On $P_n$ we have $(\sigma\vee n)\wedge\varsigma=n$ and
$(\sigma\vee(n+1))\wedge\varsigma=n+1$, and the summation range gains exactly
the top index $n$, contributing the single term $u(S_n)$. Also $n<\varsigma\le\tau$
by \ref{it:slsb-order}, so \ref{it:slsb-step} gives $S_{n+1}=S_n+\Xi_n$. By
\ref{it:slsb-pm} and \ref{it:slsb-Xi} there are two subcases: if $\xi_n=1$ then
$\Xi_n=1$ and $S_{n+1}=S_n+1$, and
\[
  h(S_n+1)=\frac{h(S_n+1)+h(S_n-1)}{2}
    +\frac{h(S_n+1)-h(S_n-1)}{2}\cdot1 ;
\]
if $\xi_n=-1$ then $\Xi_n=-1$ and $S_{n+1}=S_n-1$, and
\[
  h(S_n-1)=\frac{h(S_n+1)+h(S_n-1)}{2}
    +\frac{h(S_n+1)-h(S_n-1)}{2}\cdot(-1) .
\]
In both subcases
$h(S_{n+1})=\frac{h(S_n+1)+h(S_n-1)}{2}+c_n\xi_n$, whence
$Y_{n+1}-Y_n=\iota(h(S_{n+1})-h(S_n)+u(S_n))=\iota D_n+\iota c_n\xi_n$. This
proves \eqref{eq:slsb-split}.

\emph{The martingale term integrates to zero.} The factor $\iota c_n$ is
$\mathcal G_n$-measurable: $\iota\mathbf1_{\{\sigma\le n\}}$ is, because
$A\cap\{\sigma\le n\}\in\mathcal G_n$ by hypothesis; $\mathbf1_{\{n<\varsigma\}}$
is, because $\{n<\varsigma\}$ is the complement of
$\{\varsigma\le n\}\in\mathcal G_n$ by \ref{it:slsb-stG}; and
$\omega\mapsto h(S_n(\omega)\pm1)$ is $\mathcal G_n$-measurable by
\ref{it:slsb-meas} applied to $g=h(\cdot\pm1)$ with $k=n$. It is bounded in
absolute value by $(C_h^++C_h^-)/2$, using $|\iota|\le1$ and the bounds fixed at
the start. The variable $\xi_n$ is bounded by $1$ in absolute value by
\ref{it:slsb-pm}, hence integrable, and
$\mathbb E[\xi_n\mid\mathcal G_n]=0$ almost everywhere by \ref{it:slsb-xicond}.
The pull-out property of the conditional expectation for a bounded
$\mathcal G_n$-measurable factor therefore gives
$\mathbb E[\iota c_n\xi_n\mid\mathcal G_n]=\iota c_n\,
\mathbb E[\xi_n\mid\mathcal G_n]=0$ almost everywhere, and since
$\iota c_n\xi_n$ is integrable (bounded and measurable, $\mu$ being a
probability measure),
\begin{equation}\label{eq:slsb-mart}
  \int\iota\,c_n\,\xi_n\,d\mu
   =\int\mathbb E[\iota c_n\xi_n\mid\mathcal G_n]\,d\mu=0 .
\end{equation}

\emph{The drift term is nonpositive.} Pointwise, $\iota D_n\le0$. Indeed off
$P_n$ we have $D_n=0$; and on $P_n$ hypothesis \ref{it:slsb-win} gives
$a<S_n<b$, so \eqref{eq:slsb-super} at $i=S_n$ reads
$\frac{h(S_n+1)+h(S_n-1)}{2}-h(S_n)\le-u(S_n)$, i.e.
$\frac{h(S_n+1)+h(S_n-1)}{2}-h(S_n)+u(S_n)\le0$; multiplying by
$\iota\ge0$ preserves the inequality. The function $\iota D_n$ is measurable and
bounded, hence integrable, so by monotonicity of the integral
\begin{equation}\label{eq:slsb-drift}
  \int\iota\,D_n\,d\mu\le0 .
\end{equation}

\emph{The decrease.} Both $\iota D_n$ and $\iota c_n\xi_n$ are integrable, so by
additivity of the integral, \eqref{eq:slsb-split}, \eqref{eq:slsb-drift} and
\eqref{eq:slsb-mart},
\[
  \int Y_{n+1}\,d\mu-\int Y_n\,d\mu
   =\int\iota D_n\,d\mu+\int\iota c_n\xi_n\,d\mu\le0 ,
\]
that is, $\int Y_{n+1}\,d\mu\le\int Y_n\,d\mu$ for every $n$.

\emph{Conclusion.} We prove \eqref{eq:slsb-bound} by induction on $n$. For
$n=0$ it is the equality \eqref{eq:slsb-Y0}. If it holds for $n$, then
$\int Y_{n+1}\,d\mu\le\int Y_n\,d\mu\le\int\iota\,h(S_\sigma)\,d\mu$ by the
previous paragraph and the inductive hypothesis. This proves
\eqref{eq:slsb-bound} for every $n$.
\end{proof}
